\documentclass[11pt]{amsart}
\usepackage{a4wide,amsfonts,amssymb,stmaryrd,amscd,amsmath,latexsym,amsbsy}
%\documentstyle{article}
\theoremstyle{plain}


% %\usepackage{showframe}
% \newcommand*{\SignatureAndDate}[1]{%
%    \par\noindent\makebox[2.5in]{\hrulefill} \hfill\makebox[2.0in]{\hrulefill}%
%    \par\noindent\makebox[2.5in][l]{#1}      \hfill\makebox[2.0in][l]{Date}%
%}%



\newtheorem{theorem}{Theorem}[section]
\newtheorem{thm}{Theorem}[section]
\newtheorem*{thrm}{Theorem}
\newtheorem{prop}[theorem]{Proposition}
\newtheorem{pro}[theorem]{Proposition}
%\newtheorem{defi}{Definition}
\newtheorem{lem}[theorem]{Lemma}
\newtheorem{cor}[theorem]{Corollary}
\newtheorem{eg}{{Example}}
\newtheorem{ex}{Exercise}
\newtheorem*{sol}{Solution}
\newtheorem*{LRrule}{Littlewood-Richardson rule}
\newtheorem*{conj}{Conjecture}
%\redefine{thebibliography}{Publications}
\theoremstyle{definition}
\newtheorem*{defi}{Definition}
\theoremstyle{remark}
\newtheorem{rema}{Remark}[section]
\newcommand{\CC}{\mathbb C}
\newcommand{\ZZ}{\mathbb Z}
\newcommand{\solu}[1]{\begin{sol}{\bf (\ref{#1})}}
\newcommand{\dd}{\partial}
\newcommand{\gotg}{\mathfrak g}
\newcommand{\gotk}{\mathfrak k}
\pagestyle{plain}
\begin{document}
\section*{\bf{\huge Alexei Oblomkov}}
\bigskip

%{\bf Birth date}: 1 December 1976.

%\medskip {\bf Birth place}: Ul'yanovsk, Russia.

\begin{tabbing}
Address:\quad\quad\quad\;\= 1234H, LGRT Tower\\
                      \>University of Massachusetts\\
\>  Amherst, MA 01003-9305.\\

E-mail:\quad\quad\quad\quad oblomkov@math.umass.edu\\

Citizenship:\> Russia, US Permanent Resident since 2010, US citizen since 2014.
\end{tabbing}

%\section*{\Large\underline{Research interests}}

%%\medskip

%%\begin{tabbing}

%\noindent Representation theory, algebraic geometry
%and mathematical physics.




\section*{\Large\underline{Education}}
\begin{tabbing}
2005\quad\quad\quad\quad\quad\= Massachusetts Institute of
Technology   \quad\quad\,\,\:\:\:\:\= Ph.D. in
mathematics\\
\>\quad\quad  \>  Advisor: Pavel Etingof\\
%\>\quad \>  Title: Double Affine Hecke\\
%\>\> Algebras and Noncommutative\\
%\>\> Geometry\\ 
1999-2001\> Independent University of
Moscow\> Graduate Study\\
1999-2001\>Moscow State University\> Graduate Study\\

1998\> Independent University of Moscow\>MA in Pure Mathematics\\
1998\> Moscow State University\>MA in Pure Mathematics
\end{tabbing}


\section*{\Large\underline{Employment}}

\begin{tabbing}

2020-\quad\quad\quad\quad\;\;\;\= University of Massachusetts, Amherst\quad\quad\;\;\;  Professor\\
  
2018 Spring\quad\quad  MSRI \quad\quad\quad\quad\quad\quad\quad\quad\quad\quad\quad\quad\quad\quad\quad\quad\quad Research Professor\\
2015-2019\quad\quad\;\;\;\= University of Massachusetts, Amherst\quad\quad\;\;\; Associate Professor\\
2009-2015\quad\quad\;\;\;\= University of Massachusetts, Amherst\quad\quad\;\;\; Assistant Professor\\
2008-2009\quad\quad\;\;\;\= Princeton University\quad\quad\quad\quad\quad\quad\quad
\quad\quad\quad\;\; Visiting Assistant
Professor\\
 2006-2008\quad\quad\;\;\:\=        Princeton
University\quad\quad\quad\quad\quad\quad\quad\quad\quad\quad\quad\=
Instructor\\
2005-2006 \>Institute for Advanced Study\> Postdoctoral Fellow
\end{tabbing}






%\section*{\Large\underline{Research interests}}

%\medskip

%\begin{tabbing}

%Representation theory:\quad\quad\= Special functions and
%orthogonal polynomials,
%          Hecke algebras\\
% Algebraic geometry:\> Gromov-Witten and Donaldson-Thomas
% invariants,\\
%         \> noncommutative geometry, projective geometry\\
% Mathematical physics:\>
%          Inverse scattering, integrable systems
%\end{tabbing}

\section*{\Large\underline{Grants}}
\begin{tabbing}
  2022-2025\qquad\quad\= NSF DMS-2200798,\quad\quad
\quad\quad\quad\quad\quad\quad\quad\quad\quad\;\$217,038.00, sole PI\\
 2018-2021\qquad\quad\=NSF FRG DMS-1760373, \quad\quad\quad\quad \quad\quad\quad\quad \;  \$300,000,  for UMass\\ 
2015-2016\qquad\quad\=NSF DMS-1546311, Workshop funding\quad\quad\quad \$25,000, sole co-PI\\
2014-2019\qquad\quad\=NSF CAREER DMS-1352398\quad\quad\quad\quad\quad\quad\;\;\;\$420,000.00, sole PI\\
2010-2013\qquad\quad\=NSF DMS-1001609\quad\quad
\quad\quad\quad\quad\quad\quad\quad\quad\quad\;\;\$123,910.00, sole PI\\
2007-2010\qquad\quad\=NSF
DMS-0701367\quad\quad\quad\quad\quad\quad\quad\quad\quad\quad\quad\;\;\$111,300.00, sole PI\\
%2005-2006\> NSF Division of Mathematical Sciences Fellowship, DMS-0111298 \\
%2003-2006\> Civilian Research and Development Foundation Grant,
%RMI-2545-MO-03 \\
%1997-2000\>INTAS Student Stipend and Travel grant
\end{tabbing}

\section*{\Large\underline{Awards}}
\begin{tabbing}
2018  \quad\quad\quad\quad\quad                  Simons Fellow in Mathematics for Fall 2018 \\
2014-2019\quad\quad\quad NSF CAREER Award\\
2010-2012\quad\quad\quad Sloan Fellowship\\
2004\quad\quad\quad\quad\quad\;\= Charles and Jennifer Johnson
Prize, MIT departmental award
 for the best \\ \>paper submitted by a graduate student\\
2004, 2003\> Rogers Family Prize, MIT departmental award   for the
             best mentor-mentee\\ \>team in Summer Program for
             Undergraduate research\\
%1996-1997\>  Lomonosov Stipend, Moscow  State University award for
%high academic \\ \>achievement
\end{tabbing}




\section*{\Large\underline{Publications}}

\noindent\hangindent=0.8 cm 1.  {\it Integrability of some quantum
systems associated with the root system $B_2$,} (Russian) Vestnik
Moskov. Univ. Ser. I Mat. Mekh.  72 (1999), no. 2, 6--9;
        translation in Moscow Univ. Math. Bull. 54 (1999), no. 2, 5--8.

\noindent\hangindent=0.8 cm 2. {\it Monodromy-free Schrodinger
operators with quadratically increasing potential,} (Russian)
Teoret. Mat. Fiz. 121 (1999), no. 3, 374--386; translation in
Theoret. and Math. Phys. 121 (1999), no. 3, 1574--1584.

\noindent\hangindent=0.8 cm 3. (with F. L. Zak, A. V. Inshakov, S.
M. L'vovski) {\it On congruences of lines of order one in
$\mathbb{P}^3$,} preprint (1999), available at
http://www.mccme.ru/ium/postscript/ s99/notes.

\noindent\hangindent=0.8 cm 4. (with O.  Chalykh), {\it Harmonic
oscillator and Darboux transformations in many dimensions,} Phys.
Lett. A 267 (2000), no. 4, 256--264.

\noindent\hangindent=0.8 cm 5. (with  A.  Penskoi) {\it
Two-dimensional algebro-geometric difference operators,} J. Phys.
A 33 (2000), no. 50, 9255--9264.

\noindent\hangindent=0.8 cm 6. {\it Difference operators on
two-dimensional regular lattices,} (Russian) Teoret. Mat. Fiz. 127
(2001), no. 1, 34--46; translation in Theoret. and Math. Phys. 127
(2001), no. 1, 435--445.

\noindent\hangindent=0.8 cm 7. {\it On the spectral properties of
two classes of difference periodic operators,} (Russian) Mat. Sb.
193 (2002), no. 4, 87--112; translation in  Sb. Math.  193 (2002),
no. 3-4, 559--584.

\noindent\hangindent=0.8 cm 8. {\it The isoenergy spectral problem
for multidimensional difference operators,} (Russian) Funktsional.
Anal. i Prilozhen. 36 (2002), no. 2, 45--61;  translation in
Funct. Anal. Appl.  36  (2002),  no. 2, 120--133.

\noindent\hangindent=0.8 cm 9. (with O. Chalykh,   P. Etingof)
{\it  Generalized Lame operators,} math.QA/0212029, Communications
in Math. Physics. 239 (2003), no. 1-2, 115--153.

\noindent\hangindent=0.8 cm 10. {\it Heckman-Opdam's Jacobi
       polynomials for the $BC_n$ root system and
        generalized spherical functions,} math.RT/0202076,
        Adv. Math.,  186  (2004),  no. 1, 153--180.

\noindent\hangindent=0.8 cm 11. {\it The double affine Hecke
algebras and Calogero-Moser spaces,} math.RT/0303190,  Represent.
Theory 8 (2004), 243--266.

\noindent\hangindent=0.8 cm 12. {\it Double affine Hecke algebras of rank 1 and
affine cubic surfaces,} math.RT/0306393,  IMRN,  (2004),  no. 18,
877--912.

\noindent\hangindent=0.8 cm 13. (with J. V. Stokman)  {\it Vector valued spherical
function and Macdonald-Koornwinder polynomials,}  math.QA/0311512,
Compositio Math.
 141, (2005), no. 5,   1310--1350.



\noindent\hangindent=0.8 cm 14. (with A. Penskoi)
 {\it Laplace transformations and spectral theory of two-dimensional
semi-discrete and discrete hyperbolic Schroedinger operators,}
math-ph/0311036, IMRN,  (2005), no. 18, 1089--1126.


\noindent\hangindent=0.8 cm 15. (with P. Etingof)
{\it Quantization and orbifold cohomology,}
 math.QA/0311005, Proceedings of the JHLM Workshop,
 Contemporary Mathematics series, (2006), no. 417,  171-182.


\noindent\hangindent=0.8 cm 16. (with P. Etingof, W. L. Gan)
{\it Generalized double affine
Hecke algebras of higher rank,}  math.QA/0504089, Crelle's Journal,
(2006), no. 600, 177--201.



\noindent\hangindent=0.8 cm 17. (with P. Etingof, E. Rains)
{\it  Generalized double affine
Hecke algebras of rank $1$ and quantized Del Pezzo surfaces,}
math.QA/0406480, Advances in Mathematics, 212
(2007), no. 10, 749--796.


\noindent\hangindent=0.8 cm 18. {\it Deformed Harish-Chandra homomorphism for the
cyclic quiver,} math.RT/0504395, Mathematical Research Letters,
33 (2007), Issue 3, 359--372.

\noindent\hangindent=0.8 cm 19. (with P. Etingof, W. L. Gan, V. Ginzburg)
 {\it Ha\-rish-Cha\-n\-dra ho\-mo\-mor\-phisms and sym\-plec\-tic ref\-lec\-tion
algebras for wre\-ath-pro\-ducts,} math.RT/0511489, 
 Publications Mathematiques de l'IH\'ES,  No. 105  (2007), 91--155.


\noindent\hangindent=0.8 cm 20.
(with P. Etingof, S. Loktev, L. Rybnikov)
{\it A Lie-theoretic construction of spherical symplectic reflection algebras}, arXive:0809.3976,
  Transform. Groups  13  (2008),  no. 3-4, 541--556.



\noindent\hangindent=0.8 cm 21. (with E. Stoico) {\it Finite dimensional representations of
double affine Hecke algebra of rank $1$,} math.RT/0409256,
 Journal of Pure and Applied Algebra,  vol. 213  (2009),  no. 5, 766--771

\noindent\hangindent=0.8 cm 22. (with D. Maulik) {\it Quantum cohomology of
$\mathrm{Hilb}_m(A_n)$}, arXive:0802.2737, Journal of AMS, vol. 22  (2009), 1055--1091. 

\noindent\hangindent=0.8 cm 23. (with D. Maulik) {\it Donaldson-Thomas theory
of $A_n\times\mathbb{P}^1$}, arXive:0802.2739,  Compos. Math., 145 (2009), no. 5, 1249--1276.

\noindent\hangindent=0.8 cm 24. (with D. Maulik, A. Okounkov, R. Pandharipande)
{\it The Gromow-Witten/Donaldson-Thomas correspondence for toric 3-folds},
arXive:0809.3976, Invent. Math. 186 (2011), no. 2, 434--479.

\noindent\hangindent=0.8 cm 25. (with V. Shende) 
{\it The Hilbert scheme of a plane curve singularity and the HOMFLY polynomial of its link}, arxive:1003.1568,
  Duke Math. J. 161 (2012), no. 7, 1271--1303.

\noindent\hangindent=0.8 cm 26. (with V. Shende, J. Rassmusen)
{\it The Hilbert scheme of a plane curve singularity and HOMFLY homology of its link},  Geom. Topol. 22 (2018), no. 2, 645--691, arXive:1201.2115.


\noindent\hangindent=0.8 cm 27. (with E. Gorsky, J. Rassmusen) {\it On stable Khovanov homology of torus knots},    arXiv:1206.2226,  Exp. Math. 22 (2013), no. 3, 265--281.

\noindent\hangindent=0.8 cm 28. (with E. Gorsky, V. Shende, J. Rassmusen), {\it Torus knots and Rational DAHA},   arXiv:1207.4523,  Duke Math. J. 163, (2014), no. 15, 325--401. 

\noindent\hangindent=0.8 cm 29. (with Z. Yun) {\it Geometric Representations of graded and rational Cherednik algebras},   arXiv:1407.5685,
 Adv. Math. 292 (2016), 601--706.

\noindent\hangindent=0.8 cm 30. (with S. Nawata) {\it Lectures on knot homology},   arXiv:1407.5685,   Proceedings of "Workshop on Physics and Mathematics of Link Homology", at Centre de Recherches Mathematiques, Universite de Montreal,  137--177, Contemp. Math., 680, Amer. Math. Soc., Providence, RI, 2016.


\noindent\hangindent=0.8 cm 31. (with L. Rozansky) {\it Knot Homology and sheaves on the Hilbert scheme of points on the plane },
Selecta Math. (N.S.) 24 (2018), no. 3, 2351--2454,  arXiv:1608.03227.



\noindent\hangindent=0.8 cm 32. (with L. Rozansky) {\it Affine Braid group, JM elements and knot homology},
Transform. Groups 24 (2019), no. 2, 531--544. arXiv:1702.03569.


\noindent\hangindent=0.8 cm 33. (with L. Rozansky) {\it HOMFLYPT homology of Coxeter links}, Transform. Groups, 28 (2023), no. 3, 1245--1275,
arXiv:1706.00124.




\noindent\hangindent=0.8 cm 34. (with  Z. Yun) {\it The cohomology ring of certain compactified Jacobians},  arXiv:1710.05391. 


\noindent\hangindent=0.8 cm 35. (with A. Okounkov, R. Pandharipande) {\it  GW/PT descendent correspondence via vertex operators}, Communications in Mathematical Physics, arXiv:1806.00714.



\noindent\hangindent=0.8 cm 36. (with L. Rozansky)
{\it A categorification of a cyclotomic Hecke algebra}, arXiv:1801.06201.

\noindent\hangindent=0.8 cm 37. (with E. Carlsson)
{\it Affine Schubert calculus and double coinvariants}, arXiv:1801.09033.

\noindent\hangindent=0.8 cm 38. (with L. Kamenova, G. Mongardi) {\it Symplectic involutions of \(K3^{[n]}\) type and Kummer \(n\)
  type manifolds}, arXiv:1809.02810, Bull. Lond. Math. Soc. 54 (2022), no. 3, 894--909.

\noindent\hangindent=0.8 cm 39. (with L. Rozansky), {\it Categorical Chern character and braid groups},   arXiv:1811.03257 Adv. Math.437 (2024), 66 pages.

\noindent\hangindent=0.8 cm 40. (with L. Rozansky), {\it 3D TQFT and HOMFLYPT homology},  arXiv:1812.06340, Lett. Math. Phys. 113 (2023), no. 3, Paper No. 71, 62 pages.

\noindent\hangindent=0.8 cm 41. {\it EGL formula for DT/PT theory of local curves}, arXiv:1901.03014.

\noindent\hangindent=0.8 cm 42. {\it Notes on matrix factorizations and knot homology}, arXiv:1901.04052, Lecture Notes in Math., 2248
Fond. CIME/CIME Found. Subser.
Springer, Cham, 2019, 83--127.

\noindent\hangindent=0.8 cm 43. (with L. Rozansky), {\it Dualizable link homology }, arXiv:1905.06511.

\noindent\hangindent=0.8 cm 44. (with N. Bottman), {\it 
A compactification of the moduli space of marked vertical lines in $\mathbb{C}^2$, } arXiv:1910.02037.



\noindent\hangindent=0.8 cm 45. (with M. Moreira, A. Okounkov, R. Pandharipande)
{\it  Stationary Virasoro constraints for stable pairs on toric varieties}, Forum Math. Pi 10 (2022), 62 pages.


\noindent\hangindent=0.8 cm 46. (with L. Rozansky)
{\it  Matrix factorizations and \(gl(m|k)\)-quantum invariants}, arXiv:2212.02665.

\noindent\hangindent=0.8 cm 47. (with E. Gorsky, M. Mazin)
{\it  Generic curves and non-coprime Catalans}, Algebraic Geometry (2026), arXiv:2210.12569.

\noindent\hangindent=0.8 cm 48. (with E. Gorsky, O. Kivinen)
{\it The affine Springer fiber - sheaf correspondence}, {Advances in Mathematics},
 2025,
  vol. 464, arXiv:2204.00303.

\noindent\hangindent=0.8 cm 49. (with L. Kamenova, G. Mongardi )
{\it Fixed loci of symplectic automorphisms of \(K3^{[n]}\) and n-Kummer type manifolds
}, arXiv:2308.14692.

\noindent\hangindent=0.8 cm 50. (with O. Kivinen, G. Wyss )
{\it Discriminants and motivic integration},
  arXiv:2503.18449v2.



\noindent\hangindent=0.8 cm 51. (with I. Cavey, E. Gorsky, J. Turner )
{\it Newton-Okounkov bodies for nested Hilbert schemes},
arXiv:2510.07420v1.




\noindent\hangindent=0.8 cm 52. (with L. Rozansky )
{\it A categorification of representations of
                  $U_q(\mathfrak{gl}_{1|1})$},
 arXiv:2510.13039v1.














\section*{\Large\underline{Talks}}



\medskip


\noindent{\it UC Davis}, Geometry Seminar, 06/2024.

\noindent{\it CalTech}, Algebraic Geometry Seminar, 10/2024.


\noindent{\it EPFL, Switzerland}, Week-long lecture-course on TQFT, 10/2024.



\noindent {\it Les Diablerets workshop}, Workshop Categorical topological invariants, 01/2024.


\noindent{\it JHU, US}, Topology Seminar, 10/2023.

\noindent{\it AIM, US}, Workshop on knot homology, 08/2023.

\noindent{\it UC Davis, US}, Algebraic Geometry Seminar, 06/2023.

\noindent{\it Stony Brook University, US}, Algebraic Geometry Seminar, 12/2022.

\noindent{\it Yale University, US}, Representation Theory Seminar, 11/2022.

\noindent{\it MIT, US}, Representation  Theory Seminar, 11/2023.

\noindent{\it  Notre Damme University, US}, Workshop on Matrix Factorizations, 10/2023.

\noindent{\it Simons Center for Geometry and Physics}, Lectures on GW/PT correspondence, 09/2023.

\noindent{\it Harvard, US} Mathematical Physics Seminar, 10/2022.

\noindent{\it Washington University, US} Geometry Seminar, 05/2022.

\noindent{\it University of Hamburg}, Mathematical Physics Colloquium, 10/2021.

\noindent{\it University of Edinburgh, UK}, Geometric Representation Theory and Lower-dimensional topology, 06/2019.

\noindent{\it Schloss H\"unigen, Switzerland},  Helvetic Algebraic Geometry Seminar, 06/2019.

\noindent{\it ETH, Switzerland}, Algebraic Geometry Seminar, 06/2019.

\noindent{\it CalTech}, Hidden algeraic structures in topology, 04/2019.

\noindent{\it Columbia University,} Informal Mathematical Physics Seminar, 10/2018.

\noindent{\it ETH, Switzerland}, Algebraic Geometry Seminar, 10/2018.

\noindent{\it CIME, Cetraro, Italy}, Three lectures for Summer School on Geometric Representation theory and
Gauge Theory.

\noindent{\it MSRI}, Enumerative Geometry Seminar, 04/2018.

\noindent{\it UC Davis}, Geometry and  Topology Seminar, 04/2018.

\noindent{\it Perimeter Institute, Canada} Mathematical Physics Seminar talk, 10/2017.

\noindent{\it University ot Waterloo, Canada} Colloquium talk, 10/2017.

\noindent{\it University of Hamburg, Germany} String Math, Invited speaker, 07/2017.

\noindent{\it CalTech} Mathematical Physics Seminar Talk, 05/2017.

\noindent{\it MIT} Representation theory Seminar Talk, 03/2017

\noindent{\it Fields Institute, Canada}, One week workshop on Hall algebras and Enumerative invariants and Gauge Theory, Invited speaker, 11/2016.

\noindent{\it ICTP, Trieste, Italy}, One week workshop on Gauge Theory and related fields, Invited speaker, 09/2016.

\noindent{\it CRM, Montreal Canada}, One week workshop on algebraic cycles, Invited speaker, 06/2016.

\noindent{\it North-Eastern University}, Algebraic Geometry Seminar Talk, 03/2016.


\noindent{\it Texas A\&M}, Geometry and Combinatorics Seminar Talk, 10/2015.

\noindent{\it IAS/Park City Mathematics Institute}, Summer School on Geometric Representation theory, Invited speaker, 06/2015

\noindent{\it Stony Brook University}, One-week workshop on Knot homology at SCGP, Invited Speaker, 05/2015. 

\noindent{\it Banff},  One-week workshop on topological recursion, quantum curves and quantum invariants, Invited speaker, 06/2014. 

\noindent{\it North-Western University}, Conference on QFT and representation theory, Invited speaker, 05/2014.

\noindent{\it University of Geneva, Switzerland}, Conference on quantum invariants and geometry, Invited speaker, 04/2014.

\noindent{\it EPFL, Switzerland}, Three lectures on Geometry Seminar, 04/2014.

\noindent{\it North-Western University}, Topology Seminar, 03/2014.

\noindent{\it University of California, Davis}, Series of three lectures on knot invariants and algebraic geometry, 02/2014.

\noindent{\it University of Maryland}, Colloquium, 04/2014.

\noindent{\it Stanford University}, Algebraic Geometry Seminar, 09/2013.

\noindent{\it CRM, Montreal}, Two-Week workshop, Four lectures on knot homology, DAHA and singular curves 07/2014.

\noindent{\it Banff}, {One Week workshop on motivic DT invariants}, Invited speaker, 06/2013.

\noindent{\it MSRI}, {Conference on noncommutative Geometry}, Invited speaker, 05/2013.

\noindent{\it University of Michigan, Ann Arbor}, Lie Theory seminar, 04/2013.

\noindent{\it Rutgers University}, Mathematical Physics Seminar, 12/2011.

\noindent{\it University of North Carolina, Chapel Hill}, Geometry and Mathematical Physics Seminar, 12/2011

\noindent{\it University of North Carolina, Chapel Hill}, Colloquium, 12/2011.

\noindent{\it Oxford University, UK}, Series of Lectures on knot invariants and Hilbert scheme of points
on curves and DAHAs, 11/2011.

\noindent{\it Simons Center for Geometry and Physics}, Conference on Homological Invariants in Low dimensional 
Topology, 06/2011.

\noindent{\it Simons Center for Geometry and Physics}, Conference on Mirror Symmetry and Equivariant 
cohomology, 05/2011.

\noindent{\it University of Texas, Austin}, Geometry Seminar, 03/2011

\noindent{\it University of Michigan, Ann Arbor}, Mathematical Physics Seminar, 02/2011.


\noindent{\it Wall-crossing in Mathematics and Physics, Mathematics},
University of Illinois at Urbana-Champaign,
05/2010.

\noindent{\it University of Illinois, Chicago}, Algebraic Geometry Seminar,
03/2010.

\noindent{\it University of Wisconsin, Madison}, Geometry Seminar, 11/2009.

\noindent{\it University of Michigan, Ann Arbor}, Mathematical Physics Seminar,
10/2009.

\noindent{\it Boston University}, Geometry Seminar, 10/2009.

\noindent{\it AMS Special Session on TQFT}, University of Illinois at Urbana-Champaign,
04/2009.


\noindent{\it University of Chicago, Chicago}, Representation Theory Seminar, 11/2009.

\noindent{\it MIT}, Noncommutative Algebra Seminar, 05/2009.

\noindent{\it Cornell University}, Representation Theory Seminar, 04/ 2009.



\noindent{\it Stony Brook University}, Colloqium, 01/2009.

\noindent{\it University of Massachusetts, Amherst}, Special seminar, 12/2008.

\noindent{\it University of Texas, Austin}, Special seminar, December, 12/2008.

\noindent{\it University of Illinois at Urbana-Champaign}, Algebraic geometry seminar, 12/2008.

\noindent{\it University of Michigan}, Special seminar, 11/2008.




\noindent{\it Northwestern University}, Mathematical
Physics Seminar, 10/2007.

\noindent{\it  University of California, Riverside}, Lie Theory
Seminar, 12/2006.

\noindent{\it University of Virginia, Charlottesville}, Colloqium:
11/2006.

\noindent {\it University of Virginia, Charlottesville}, Algebra
Seminar, 11/2006.

\noindent{\it University of Massachusetts, Amherst},
Representation Theory Seminar, November, 12/2005.

\noindent{\it AMS Special Session on Noncommutative Algebra and
Noncommutative Birational Geometry}, University of Oregon,
11/2005.



\noindent{\it University of Washington}, Algebra Seminar,
11/2005.

\noindent{\it AMS Special Session on Algebraic Geometry and
Combinatorics}, UC, Santa Barbara, 04/2005.

\noindent{\it Rutgers University}, Lie Group Seminar, 10/2005.

\noindent{\it University of Michigan}, Algebra Seminar, 12/2004.

%\noindent{\it CalTech}, Analysis Seminar, 11/2004.

\noindent{\it MIT}, Infinite-Dimensional Algebra Seminar, 04/2004.

\noindent{\it Yale University}, Geometry, Symmetry and Physics,
12/2003.

\noindent{\it University of Massachusetts}, Amherst,
Representation Theory
          Seminar, 12/2003.

\noindent {\it Cornell University}, Lie Groups Seminar, 01/2003.



\noindent{\it MIT}, Infinite-Dimensional Algebra Seminar,
11/2002.

\noindent{\it University of Amsterdam}, Lectures on special
functions, 07/2002.

\medskip

%\newpage

\section*{ {\Large\underline{Teaching}}}
\noindent University of Massachusetts
\begin{tabbing}

Summer 2025: \quad\= REU program\\
\>Mentored four students \\
  
  
  Spring 2025:\quad\quad\= Graduate Algera, (Math 613)\\
\>Taught lectures \\

Spring 2024:\quad\quad\= Introduction into Linear Algebra, (Math 235)\\
\>Taught lectures \\


  

  Fall 2024:\quad\quad\= Graduate Algebra, (Math 612)\\
\>Taught lectures \\

  

Spring 2024:\quad\quad\= Graduate Complex Analyis, (Math 621)\\
\>Taught lectures \\

Spring 2024:\quad\quad\= Introduction into Linear Algebra, (Math 235)\\
\>Taught lectures \\



  
Fall 2023:\quad\quad\= Introduction to Complex Analyis, (Math 421)\\
\>Taught lectures \\

  
Spring 2023:\quad\quad\= Introduction to Linear Algebra, (Math 235)\\
\>Taught lectures \\




  
Fall 2022:\quad\quad\quad\= Representation Theory, (Math 680)\\
\>Taught lectures \\

  
Fall 2022:\quad\quad\quad\= Putnam Seminar, (Math 490)\\
\>Taught lectures \\


Spring 2022:\quad\quad\= Introduction to Complex Analysis, (Math 421)\\
\>Taught lectures \\


Fall 2021:\quad\quad\quad\= Multi-variable Calculus, (Math 233)\\
\>Taught lectures \\

  
Spring 2021:\quad\quad\= Graduate Topology II, (Math 672)\\
\>Taught lectures \\

  
Fall 2020:\quad\quad\quad\= Multi-variable Calculus, (Math 233)\\
\>Taught lectures \\

  
Fall 2020:\quad\quad\quad\= Graduate Topology I, (Math 671)\\
\>Taught lectures \\
  
  
Spring 2020:\quad\quad\= Introduction to Linear Algebra, (Math 235)\\
\>Taught lectures \\


  
Spring 2019:\quad\quad\= Graduate Topology II, (Math 672)\\
\>Taught lectures \\

  
Fall 2019:\quad\quad\quad\= Introduction to Linear Algebra, (Math 235)\\
\>Taught lectures \\

  
Fall 2019:\quad\quad\quad\= Graduate Topology I, (Math 671)\\
\>Taught lectures \\


Summer 2019:\quad\= REU program\\
\>Mentored four students \\
  
  
Spring 2019:\quad\quad\= Advanced Multi-variable calculus, (Math 425)\\
\>Taught lectures \\


Summer 2018:\quad\= REU program\\
\>Mentored four students \\

  
Fall 2017:\quad\quad\quad\= Introduction to Linear Algebra, (Math 235)\\
\>Taught lectures \\


Summer 2017:\quad\= REU program\\
\>Mentored four students \\



  
Fall 2016:\quad\quad\quad\= Introduction to Topology, (Math 671)\\
\>Taught lectures \\


Fall 2016:\quad\quad\quad\= Introduction into Linear Algebra, (Math 235)\\
\>Taught lectures \\

Summer 2016:\quad\= REU program\\
\>Mentored four students \\

Fall 2015:\quad\quad\quad\= Introduction to Linear Algebra, (Math 235)\\
\>Taught lectures \\


Fall 2015:\quad\quad\quad\= Introduction to Representation Theory, (Math 797)\\
\>Taught lectures \\

Summer 2015:\quad\= REU program\\
\>Mentored four students \\

Spring 2015:\quad\quad\= Introduction to Linear Algebra, (Math 235)\\
\>Taught lectures, course chair for 5 other sections of the course\\


Fall 2014:\quad\quad\quad\= Introduction to Topology, (Math 671)\\
\>Taught lectures \\


Fall 2014:\quad\quad\quad\= Introduction into Linear Algebra (Honors), (Math 235H)\\
\>Taught lectures \\



Fall 2013:\quad\quad\quad\= Knot Invariants, (Math 797KI)\\
\>Taught lectures \\


Fall 2013:\quad\quad\quad\= Introduction into Linear Algebra, (Math 235)\\
\>Taught lectures \\


Summer 2013:\quad\= REU program\\
\>Mentored one student \\


Fall 2012:\quad\quad\quad\= Putnam Exam, (Math 491A)\\
\>Taught lectures \\

Fall 2012:\quad\quad\quad\= Multivariable Calculus, (Math 233)\\
\>Taught lectures \\

Fall 2012:\quad\quad\quad\= Introduction into Linear Algebra, (Math 235)\\
\>Taught lectures \\

Fall 2011:\quad\quad\quad\= Introduction into Linear Algebra, (Math 235)\\
\>Taught lectures \\
Fall 2010:\quad\quad\quad\= Graduate Algebra, (Math 631)\\
\>Taught lectures \\
Spring 2010:\quad\quad\= Calculus I, (MAT 131, 2 sections)\\
 \>Taught lectures and problem sessions\\
\end{tabbing}


\noindent Princeton University:
\begin{tabbing}
Spring 2009:\quad\quad \=Multivariable Calculus (MAT
201)
\\ \>Taught  review sessions\\

Spring 2009:\quad\quad \=Junior Seminar: Representation Theory of Finite Groups
\\ \>Organize the Seminar\\



Fall 2008:\quad\quad\quad \=Advanced Multivariable Calculus (MAT
203, Course Head)
\\ \>Taught lectures, review sessions\\


Spring 2008:\quad\quad \= Multivariable Calculus (MAT
201, 2 sections)
\\ \>Taught lectures and problem sessions\\


Fall 2007:\quad\quad\quad \=Advanced Multivariable Calculus (MAT
203, 2 sections)
\\ \>Taught lectures and problem sessions\\

Spring 2007:\> Multivariable Calculus (MAT 201, 2 sections) \\
\>Taught lectures
and problem sessions\\

Fall 2006:\> Advanced Multidimensional Calculus (MAT 203, 2
sections)\\ \>Taught lectures and problem sessions \\
\end{tabbing}

\medskip \noindent MIT:

\begin{tabbing}
Summer 2004:\quad\=Summer Program in Undergraduate Research at
                   MIT\\\>
Mentored 2 undergraduate students \\
Spring 2004:\>Linear Algebra (18.06)\\ \> Taught problem sessions\\
Fall 2003:\>
          Advanced Calculus (18.01A)\\ \>
         Multidimensional Calculus (18.02)\\ \>
         Taught problem sessions\\
Summer 2004: \> Summer
Program in Undergraduate Research at MIT\\
\> Mentored 2 undergraduate students\\
Spring 2003:\> Differential Equations (18.03) \\ \> Taught problem
sessions and graded problems\\
2001-2002:\> Graded problems for 2 graduate level courses
\end{tabbing}

\medskip
%\newpage
\noindent Independent University of Moscow:
\begin{tabbing}
1999-2001:\quad\quad\quad\=  Advanced Algebra courses\\ \> Taught
problem sessions\\
%Spring 2000: \>Organized the seminar  on Projective Algebraic\\
%\>Geometry for undergraduate students.
\end{tabbing}











\section*{\Large \underline{Services}}

\noindent {\bf External Service}:

\medskip

\medskip

\noindent {\bf Refereeing:} Communications in Mathematical Physics, Contemporary
Mathematics, Transformation Groups, IMRN, JAMS, Annals of Mathematics, Journal of Algebraic Combinatorics A, Advances in Mathematics, Transactions of AMS, Compositio, Journal of Algebra, Selecta, Inventiones.
\medskip  

\noindent {\bf Refereeing Grant Proposals:} NSERC, NSA, Swiss National Science Foundation,
 the Netherlands Organisation for Scientific Research, the Chilean National Science and Technology Commission,
Newton Institute Scientific proposals, served on three NSF panels.


\medskip

\noindent{\bf Conference organization:}
PI co-organized one-week workshop in Banff (June 2014); co-organized three-day workshop in Amherst (October 2015); co-organized
three-day worshop in Amherst (July 2017); co-organize one-week AIM workshop (October 2018); co-organized one-week FRG workshops
at UC Davis (June 2019), at UMass (June 2020), at University of Oregon (2022); co-organized AIM community (2021-2022),
co-organized two-month program at Simons Center for Physics and Geometry (2023), co-organized AIM workshop (2023), co-organizer of
month-long program at Bernoulli Institute at EPFL, Lausanne (2025).

\medskip

\medskip

\noindent {\bf Internal Service:}

\medskip

\noindent Princeton University:



\noindent Coorganizer of Princeton Algebraic Geometry Seminar (Spring, 2009); served on six general exam committees; advised an undergraduate student on his senior thesis.


\medskip

\noindent University of Massachusetts:

\noindent Coorganizer of Valley Geometry Seminar and Representation Theory Seminar (2009-), Quantum Field Theory Seminar (Fall 2010);
gave   lectures on Algebraic Geometry Reading Seminar (Spring 2009: 5 lectures, Spring 2010: 2 lectures, Fall 2010: 2 lectures, Spring 2011: 2 lectures, Spring 2012: 2 lectures); 
served on Colloquium Committee ( Fall 2010, Fall 2013, Spring 2014, 2016-2017); 
Served on Advanced Graduate Algebra Exam Committee (Fall 2011-- Spring 2013); Served on Advanced Graduate Topology Exam Committee Committee (Fall 2014-- Spring 2023);
Member of Faculty Search Committee 
(Fall 2011, Spring 2012, Fall 2012, Spring 2013, Spring 2019); Chair the Search Committee for Visiting Assistant Professor (2018-)
Undergraduate Advisor (2010-), Served on Undergraduate Affairs Committee (Fall 2017); Member of Climate Committee (2023);
Trained UMass Putnam Team (2022, 2015, 2016). 

% \noindent During the summers 2003 and 2004 PI participated Summer
% Program in Undergraduate Research (SPUR) as a mentor of two (2003)
% and one  undergraduate students. SPUR is a six-week program of
% full-time research experience for MIT undergraduates, paired with
% MIT Math graduate students as mentors, culminating in written
% papers and lectures to faculty. Both times our teams were selected
% as top teams.

\medskip

\noindent During the summer 2013 PI was supervising REU project of
UMass for one student. During summers 2025, 2019, 2018, 2017, 2016, 2015 PI supervising
REU projects for the teams of four students. 

\newpage

\section*{\Large \underline{Ph.D. advising}}

\medskip

1. My student Tom Shelly graduated in August 2016 and started his job as Visiting Assistant Professor at Mount Holyoke College. 

2. My student Toby Willson graduated in December 2015 and started his job Bloomberg. 

3. My  student Arthur Wang graduated in May 2024.

4. My current student Pranav Kalkunte works on a problem related to motivic integration.

5. My current student Chuijiao Zhang works on  Springer theory for DAHAs.

%\section*{\Large \underline{References}}

%\begin{tabbing}
%Professor P. Etingof\quad\quad\quad\quad\= MIT, USA, etingof@math.mit.edu \\
%Professor A. Okounkov\> Princeton University, USA,
%okounkov@math.princeton.edu\\
%Professor R. Pandharipande\>Princeton University, USA,
%rahulp@math.princeton.edu
%\end{tabbing}


\vspace{25mm}
\begin{tabular}{@{}p{2.5in}p{1.5in}p{2in}@{}}
\hrulefill && Date: 28 November, 2024\\
Alexei Oblomkov &&\\
University of Massachusetts, Amherst
\end{tabular}




\end{document}


%%% Local Variables:
%%% mode: latex
%%% TeX-master: t
%%% End:
